%!TEX root = ../template.tex
%%%%%%%%%%%%%%%%%%%%%%%%%%%%%%%%%%%%%%%%%%%%%%%%%%%%%%%%%%%%%%%%%%%
%% chapter1.tex
%% NOVA thesis document file
%%
%% Chapter with introduction
%%%%%%%%%%%%%%%%%%%%%%%%%%%%%%%%%%%%%%%%%%%%%%%%%%%%%%%%%%%%%%%%%%%

\typeout{NT FILE chapter1.tex}%

\chapter{Introduction}
\label{cha:introduction}


\section{Context}
\label{sec:context}

The study of Formal Languages and Automata Theory (FLAT) is integral to 
Computer Engineering and Computer Science programs. However, due to its 
abstract and mathematical nature, students often struggle to fully grasp 
the concepts and models involved. Therefore, a tool that complements the 
study of FLAT theory by being practical, user-friendly, and visually 
oriented is highly beneficial.

In recent years, the OCaml-FLAT/OFLAT tool has been developed at FCT-UNL, 
offering extensive support in this field. Despite its advanced state, 
with support for several key models, there are still many features that 
could enhance its functionality.

This dissertation aims to introduce new capabilities to OCamlFLAT/OFLAT, 
enhancing both the OCamlFLAT library and the OFLAT graphical interface.


\section{Objectives}
\label{sec:objectives}

The primary objective of this dissertation is to incorporate support for 
attribute grammars. Attribute grammars extend context-free grammars by 
allowing additional information (attributes) to be associated with 
production rules. These attributes are crucial for describing the 
semantic properties of the languages generated by the grammars.

To achieve this, the OCamlFLAT library needs to be extended to support 
attribute grammars. Additionally, graphical enhancements are required to 
enable the creation and editing of these grammars within the OFLAT 
platform.


\section{Document Structure}
\label{sec:document_structure}

This thesis is structured into multiple chapters, each covering essential aspects of the research, development, and implementation of the OCamlFLAT library and the OFLAT platform. The document follows a logical progression from theoretical background to practical applications. Below is an overview of its structure:

\begin{itemize}
    \item \textbf{Chapter 1: Introduction} – Presents the motivation behind this research, its objectives, and a general overview of the thesis.
    \item \textbf{Chapter 2: Background and Related Work} – Discusses the theoretical foundations of Formal Language and Automata Theory (FLAT), the OCaml programming language, and existing related tools. It also provides an overview of similar research efforts.
    \item \textbf{Chapter 3: System Design} – Details the design principles of the OCamlFLAT library and the OFLAT platform, describing their architecture, module organization, and main functionalities.
    \item \textbf{Chapter 4: OCamlFLAT and OFLAT Systems} – Explains the implementation of the OCamlFLAT library, its core modules, and the functionalities they provide. It also describes the OFLAT platform’s graphical interface and how it facilitates interaction with formal models.
    \item \textbf{Chapter 5: Declarative Code in OCaml} – Explores the advantages of declarative programming in OCaml, particularly in the context of formal models. It also discusses JavaScript interoperability using \texttt{js\_of\_ocaml}, the representation of attributes, and visualization with Cytoscape.
\end{itemize}



